\section{Introduction}

Diabetic retinopathy is a leading cause of preventable blindness among working-age adults \citep{morgan2021burden}.
Referral grading is commonly performed by trained human graders, which is time-consuming and subject to inter-observer variability.
Automated screening has therefore attracted increasing attention, especially where specialist graders are scarce.

Previous studies have applied convolutional networks to referral grading and reported sensitivities above 90\% \citep{example2022deep}.
Rivera et al. improved lesion detection by adding a multi-scale attention module \citep{rivera2023attention}.
Despite these advances, explaining the reason for referral remains an open problem.
Most of these models are trained only with image-level labels, so they may rely on camera artefacts rather than the lesions themselves \citep{sample2024shortcut}.

In this study, we aim to make referral decisions traceable to the lesions that justify them.
To this end, we designed a two-stage network that grades referral urgency only from the lesions it detects.
